\section{Question 2: Capacity-valid scheduling}
\label{sec:r2q2}
\subsection{Contract and Algorithm 1}
Each GPU model is simulated as its own pool. A job starts at $s_j\ge a_j$, completes at $f_j=s_j+p_j$, and all $W_j$ workers start together. At a proposed start time, a placement $x_{ji}$ is admitted only if the \emph{pre-start} residual capacities satisfy
\begin{equation}
 g_jx_{ji}\le G_i^{free}(t),\qquad c_jx_{ji}\le C_i^{free}(t),\qquad \sum_i x_{ji}=W_j.
\end{equation}
The resulting full-time model is instead the active-job constraint $\sum_{j:s_j\le t<f_j}g_jx_{ji}\le G_i^{cap}$ and its CPU analogue. Thus a single job's check against residual capacity is not written as a whole-cluster dynamic constraint. Durations are exact integer seconds; gang admission is atomic, non-preemptive and non-migrating.

\takeaway{Algorithm 1 (SPT-HOL)}{For each model pool, insert only arrived jobs into a shortest-processing-time heap keyed by $(p_j,a_j,\mathrm{id}_j)$. Examine the heap head. If it is infeasible on the current residual nodes, stop dispatching this pool at this timestamp (head-of-line blocking). If feasible, scan nodes in the selected placement order, verify every worker slot's GPU and CPU residual, then admit all workers atomically and schedule one completion event. For GPU-series-1 only, the main method uses a CPU/GPU-ratio node order before the same full slot check; every other pool uses ascending-ID first fit. Advance to the next arrival or completion and repeat.}

\method{SPT-FILL} uses the same key but may inspect later feasible jobs under the fixed $K=200$ bypass budget. A heap insertion or removal costs $O(\log n)$; a dispatch also scans candidate nodes and worker slots, and ratio placement orders the candidate nodes by shape. Multiple starts at one timestamp can repeat these scans, so an $O(\log n)$ event claim would be incomplete. Saved schedules were independently checked by route, physical-capacity and metric validators.

The simultaneous-release mechanism has a precise necessary condition. If $f^-$ is the number of free worker slots immediately before an event, $\Delta f$ is the number released by jobs finishing at that event, and $k$ is the next job width, admission requires
\begin{equation}
 f^-+\Delta f\ge k.
\end{equation}
Only with a continuously feasible single-worker backlog and zero free slots after the preceding dispatch does this reduce to $\Delta f\ge k$. A finite $K$ candidate budget can leave $f^->0$, so the reduction is not generally valid. Multiple-worker jobs may wait when free slots cannot accumulate into a gang-sized release; HOL can allow free slots to accumulate while the head is blocked. The formal grouped results in Table~\ref{tab:r2strata} are consistent with this mechanism; it is not the unique proved cause of HOL's result or a starvation-free guarantee.

\begin{table}[H]\centering
\caption{Six-pool SPT-HOL results. GS1 uses ratio; other pools use FF-id. Waiting uses complete drain; allocation and completion use $H=7{,}776{,}000$ s.}
\label{tab:r2pools}
\begin{tabular}{lrrrrrr}
\toprule
Pool & Jobs & GPUs & \shortstack{Mean\\wait (h)} & \shortstack{Median/P99\\wait (h)} & \shortstack{GPU/CPU\\allocation} & \shortstack{90-day\\done}\\
\midrule
A10 & 111,284 & 2,494 & 32.906 & 0.023/1542.181 & 99.922\%/17.927\% & 107,002\\
A100 & 80,877 & 3,456 & 0.000 & 0.000/0.000 & 62.755\%/42.032\% & 80,101\\
GS1 & 9,643 & 1,558 & 0.000 & 0.000/0.000 & 60.867\%/34.814\% & 8,221\\
GS2 & 13,633 & 976 & 0.016 & 0.000/0.045 & 56.931\%/37.424\% & 13,606\\
H800 & 4,716 & 1,752 & 0.016 & 0.000/0.000 & 55.816\%/52.193\% & 4,705\\
A800 & 3,959 & 176 & 0.115 & 0.000/4.390 & 57.311\%/49.493\% & 3,943\\
\bottomrule
\end{tabular}
\smallskip\par A100 and A800 denote the SXM4-80GB models; GS1 and GS2 denote GPU-series-1 and GPU-series-2. Allocation integrates $[0,H)$; a completion at $H$ is included.
\end{table}


The source and configuration mapping for Table~\ref{tab:r2pools} is in the supplement. Exact per-job comparisons confirm that this GS1 placement change leaves A10 timing unchanged.

\subsection{A count lower bound and exact accounting}
For the complete A10 set let $C=2,494$, $G_j=W_jg_j$, $Z(t)=\{j:a_j\le t<a_j+p_j\}$, and let $q(t)$ be the minimum number of jobs removed from $Z(t)$ so that the remaining GPU width is at most $C$. For a schedule, let $S(t)=\{j:s_j\le t<f_j\}$ and $A(t)=S(t)\cap Z(t)$. At any $t$, at least $q(t)$ jobs in $Z(t)$ cannot be simultaneously active, hence $|Z(t)\setminus A(t)|\ge q(t)$. Integrating and using the waiting identity gives
\begin{equation}
\sum_j\min(\omega_j,p_j)=\int |Z(t)\setminus A(t)|\,dt\ge\int q(t)\,dt,
\qquad \sum_j\omega_j\ge\int q(t)\,dt,
\end{equation}
so $L_{count}=(3600N)^{-1}\int q(t)dt=22.3179305700$ h is a lower bound on equal-weight mean waiting. For the observed schedule,
\begin{equation}
\omega_j=\min(\omega_j,p_j)+\max(\omega_j-p_j,0).
\end{equation}
The exact full-A10 accounting, shown with display rounding, is
\begin{equation}
32.9064540350 = 22.3179305700 + 8.3361103922 + 2.2524130728\ \mathrm{h}.
\end{equation}
The first term is the count lower bound, the second is the within-window accounting gap, and the third is waiting beyond a job's own duration. The displayed sum is rounded to the shown precision; these are accounting terms rather than causal shares. The conservative certificates $\bar\omega/L$ are 1.4744 for ordinary job mean and 1.2090 for GPU-work-weighted mean (the latter compares 677.689933 h with the 560.521368 h weighted bound); they are for this fixed instance and different objectives.

\subsection{A10 comparison and worker strata}
For the same complete 111,284-job A10 population and fixed native GPU-work weights $H_j=W_jg_jp_j$, the three formal schedules give:
\begin{table}[H]\centering
\caption{Complete A10 waiting comparison using the same 111,284 jobs and fixed native GPU-work weights.}
\label{tab:r2a10}
\begin{tabular}{lrrrr}
\toprule
Policy & Mean (h) & P99 (h) & Max (h) & \shortstack{GPU-work-weighted\\mean (h)}\\
\midrule
FCFS & 900.129 & 1419.616 & 1423.254 & 619.352\\
SPT-HOL & 32.906 & 1542.181 & 3017.288 & 677.690\\
SPT-FILL & 34.967 & 1575.462 & 3015.311 & 677.496\\
\bottomrule
\end{tabular}

\end{table}


FCFS has lower P99 and maximum waiting than SPT-HOL, but its ordinary mean is \FCFSratio{} times higher; with fixed GPU-work weights, FCFS instead has the lower mean. The supplement identifies the complete formal schedules used in Tables~\ref{tab:r2a10} and~\ref{tab:r2strata}.

SPT-HOL is the average-wait choice on this fixed trace: it lowers ordinary mean waiting by 2.061 h relative to SPT-FILL and lowers P99 by 33.280 h, but its maximum is 1.978 h higher. The management trade-off becomes clearer under the fixed GPU-work objective: the ranking reverses, with SPT-FILL at 677.496 h versus SPT-HOL at 677.690 h. Thus ``lower average but a higher extreme'' and ``different ranking after weighting'' are both observed facts, not dominance claims.

\begin{table}[H]\centering
\caption{HOL/FILL worker strata: waiting in hours; paired difference is HOL minus FILL.}
\label{tab:r2strata}
\begin{tabular}{lrrrrrrr}
\toprule
Stratum & Jobs & \shortstack{HOL\\mean} & \shortstack{FILL\\mean} & \shortstack{HOL\\P99} & \shortstack{FILL\\P99} & \shortstack{Paired\\mean} & \shortstack{Improve / worsen\\/ same (jobs)}\\
\midrule
$W=1$ & 110,497 & 33.125 & 32.805 & 1553.319 & 1539.480 & $+0.320$ & \shortstack{54,489 / 50,566\\ / 5,442}\\
$W=2$--8 & 637 & 1.529 & 186.767 & 15.603$^\dagger$ & 1952.714$^\dagger$ & $-185.237$ & \shortstack{628 / 8\\ / 1}\\
$W\ge9$ & 150 & 4.997 & 982.722 & 36.186$^\dagger$ & 2907.896$^\dagger$ & $-977.725$ & \shortstack{149 / 1\\ / 0}\\
\bottomrule
\end{tabular}
\smallskip\par $^\dagger$ Small strata; P99 is unstable. The groups use the same complete A10 population.
\end{table}


The strata show why a pooled mean can hide policy preference: HOL is much better for multi-worker jobs, whereas FILL has slightly lower mean and P99 for the dominant single-worker stratum. The paired full-trace totals are 55,266 improved, 50,575 worsened and 5,443 unchanged for HOL minus FILL, with mean difference $-2.060576$ h. These values compare the same A10 IDs.

For the fixed GS1 pool, ratio placement is a constructive zero-wait schedule: all 9,643/9,643 jobs have zero waiting and mean waiting 0 h, while the FF-id baselines are 0.029188 h (FCFS) and 0.434277 h (SPT-HOL). Since mean waiting is nonnegative, 0 is the optimum for this fixed GS1 instance. The ratio rule is not generalized to other heterogeneous pools.

\subsection{Second-round policy evidence}
SPT-HEAD-TIME is a bounded policy with a precise implementation definition. When the SPT head is blocked, the code builds a reservation for the first future release event at which the head has a node-level placement witness. While that reservation is active, a backfill candidate is admitted only when it is currently feasible and its finish time is no later than the promised release; a candidate that crosses the promise is rejected by the TIME rule. The $K=200$ budget counts candidate attempts only after the blocked head has started a bypass scan. Under the shared-timer event mode, timer thresholds at each arrival plus $\theta$ create additional dispatch events; for this policy they do not mark jobs urgent. These statements follow the implementation recorded in the supplement, which gives the exact source lines and saved summary.

On the full A10 trace, SPT-HEAD-TIME ($\theta=168$ h, $K=200$, shared timer) changes mean waiting from 32.906454 to 32.923817 h, P99 from 1542.181042 to 1541.586592 h, and maximum from 3017.288056 to 3017.727778 h; its higher cost is not hidden by a single tail statistic. The K/timer and HEAD-TIME results support a bounded mechanism interpretation, not starvation freedom or a general performance theorem. The urgent policies remain prefix-only and are not promoted to the complete-trace recommendation.
